% 本书（《大学物理》第 3 册，第一学年）照片与 AI 插图的图片来源。
% 机器可读的许可记录位于 images/book3/CREDITS.md；两个文件须保持同步。
\chapter*{图片来源}

本书中的插图均为原创。照片依其各自的自由许可协议使用，谨向作者致谢：

\begin{itemize}
  \item \emph{基布尔秤（NIST-4）}（第 1 章）：J.~L.~Lee，NIST 提供，公有领域。
  \item \emph{叶凯士 40 英寸折射望远镜}（第 4 章）：Kb9vrg，Wikimedia Commons，公有领域。
  \item \emph{月球正面}（第 4 章）：NASA/GSFC/Arizona State University (LRO)，公有领域。
  \item \emph{吉他琴弦}（第 5 章）：Acagastya，Wikimedia Commons，CC0。
  \item \emph{示波器，正弦波与方波}（第 8 章）：Xato，Wikimedia Commons，公有领域。
  \item \emph{过山车环形轨道}（第 11 章）：Jeremy Thompson，Wikimedia Commons，CC~BY~2.0。
  \item \emph{国际空间站}（第 12 章）：NASA，公有领域。
  \item \emph{GPS 卫星}（第 16 章）：NASA，公有领域。
  \item \emph{伯克利 60 英寸回旋加速器}（第 17 章）：US Department of Energy，公有领域。
  \item \emph{傅科摆，先贤祠}（第 18 章）：Olga Khomitsevich，Wikimedia Commons，CC~BY~2.0。
  \item \emph{潜艇浮出水面}（第 21 章）：US Navy，公有领域。
  \item \emph{闪电}（第 26 章）：Mark Coldren，Wikimedia Commons，CC0。
  \item \emph{密立根油滴仪}（第 27 章）：Daderot，Wikimedia Commons，CC0。
  \item \emph{指南针}（第 28 章）：Evan-Amos，Wikimedia Commons，公有领域。
  \item \emph{高压输电线路}（第 28 章）：Stefan Andrej Shambora，Wikimedia Commons，CC~BY~2.0。
  \item \emph{核磁共振扫描仪}（第 29 章）：Jan Ainali，Wikimedia Commons，CC~BY~3.0。
  \item \emph{氢原子发射光谱}（第 30 章）：Merikanto and Adrignola，Wikimedia Commons，CC0。
\end{itemize}

每张照片的完整来源链接与许可协议网址，均列于本书代码库中的 \texttt{images/book3/CREDITS.md}。

\bigskip
除照片外，第 1--4 章、第 6--9 章、第 12--15 章、第 19 章、第 21--25 章及第 27--30 章中共有二十四幅图为 AI 生成插图，由 OpenAI 图像生成工具依本书编辑撰写的提示词制作，并在收录前已审核其物理内容的准确性。每幅生成图像的完整提示词均记录于代码库中的 \texttt{images/book3/ai/PROMPTS.md}。
\clearpage
